% QAOA path selection for network traffic routing -- formulation and solution.
% Qubit hackathon (Classiq), pipeline steps 3-4.
% Build: pdflatex qaoa_routing.tex (twice for cross-references).
\documentclass[11pt]{article}

\usepackage[margin=1in]{geometry}
\usepackage{amsmath, amssymb}
\usepackage{booktabs}
\usepackage{microtype}
\usepackage[colorlinks=true, linkcolor=blue, citecolor=blue, urlcolor=blue]{hyperref}

\newcommand{\onehot}{\mathrm{1hot}}
\newcommand{\lat}{\mathrm{lat}}
\newcommand{\cng}{\mathrm{cong}}
\newcommand{\switch}{\mathrm{switch}}
\newcommand{\used}{\mathrm{used}}
\newcommand{\Uncap}{\mathrm{Uncap}}

\title{Quantum Path Selection for Network Traffic Routing:\\
a QAOA Formulation and Solution}
\author{Qubit Hackathon team}
\date{\today}

\begin{document}
\maketitle

\begin{abstract}
We address the AT\&T network-traffic-routing challenge with a hybrid
quantum--classical pipeline: classical preprocessing enumerates a small set of
feasible candidate paths per traffic demand, a quantum optimizer (QAOA on the
Classiq platform) selects one path per demand by minimizing a QUBO cost
Hamiltonian that trades off latency against congestion, and classical
post-processing validates feasibility and computes operational KPIs. This
document states the QUBO formulation, the QAOA solution method as implemented,
and a worked example on a 9-qubit toy instance, including the
Fortz--Thorup congestion cost \cite{fortz2000} used both as an optional
in-Hamiltonian congestion profile and as an instance-independent KPI.
\end{abstract}

\section{Problem and hybrid pipeline}
\label{sec:problem}

A communication network is a directed graph $G=(V,E)$ whose links $e \in E$
have finite capacity $c_e > 0$ and latency $\ell_e \ge 0$. Traffic consists of
$K$ demands: demand $k$ must route bandwidth $b_k > 0$ from source $s_k$ to
target $t_k$. The goal is to select routes that use resources efficiently
while avoiding congestion and holding service quality --- the challenge KPIs
are congestion relief on the most loaded links, low average latency, high
demand satisfaction, few route changes against the current state, and zero
capacity violations.

The solution is a five-step hybrid pipeline; this document covers the
combinatorial core (steps 3--4):

\begin{enumerate}
  \item \textbf{Input} the network graph, link capacities/latencies, and the
        demand matrix.
  \item \textbf{Classical path generation:} for each demand $k$, compute a
        small candidate set $P_k$ of feasible paths ($K$-shortest paths,
        filtered by latency/policy constraints). The candidate sets are an
        \emph{input} to the steps below.
  \item \textbf{QUBO construction (Section~\ref{sec:qubo}):} one binary
        variable per candidate path; a cost Hamiltonian that selects exactly
        one path per demand, prefers low latency, and penalizes congestion.
  \item \textbf{Quantum optimization (Section~\ref{sec:qaoa}):} QAOA over the
        cost Hamiltonian, implemented on Classiq; a classical optimizer tunes
        the variational parameters.
  \item \textbf{Classical decode and validation
        (Section~\ref{sec:decode}):} map sampled bitstrings back to routes,
        check feasibility, and compute KPIs against a classical baseline.
\end{enumerate}

Restricting each demand to a precomputed candidate set makes this a
\emph{restricted optimal general routing} (explicit-path, MPLS-style) problem
in the sense of Fortz--Thorup \cite{fortz2000}, and keeps the qubit count at
$N = \sum_k |P_k|$ --- independent of the number of links.

\section{QUBO formulation}
\label{sec:qubo}

\subsection{Encoding}

The decision variables are
\begin{equation}
  x_{k,p} \in \{0,1\}, \qquad
  x_{k,p} = 1 \iff \text{demand } k \text{ routes on its candidate path }
  p \in P_k .
\end{equation}
Variables are packed into one flat register of
$N = \sum_{k=1}^{K} |P_k|$ qubits by
\begin{equation}
  i(k,p) = o_k + p, \qquad o_k = \sum_{j<k} |P_j| ,
\end{equation}
which supports differing candidate-set sizes $|P_k|$ per demand.

\subsection{Cost Hamiltonian}

The full (diagonal) cost Hamiltonian is
\begin{equation}
  \label{eq:H}
  H \;=\; \lambda_{\onehot}\, H_{\onehot}
     \;+\; \widetilde{H}_{\lat}
     \;+\; \lambda_{\cng}\, \widetilde{H}_{\cng}
     \;+\; \lambda_{\switch}\, H_{\switch},
\end{equation}
with all coefficients precomputed classically as plain floats. The same
polynomial is used symbolically to build the quantum circuit and numerically
to evaluate sampled bitstrings (Section~\ref{sec:qaoa}).

\paragraph{One path per demand (constraint penalty).}
\begin{equation}
  \label{eq:onehot}
  H_{\onehot} \;=\; \sum_{k=1}^{K} \Bigl( \sum_{p \in P_k} x_{k,p} - 1 \Bigr)^{2}.
\end{equation}
$H_{\onehot} = 0$ exactly on one-hot assignments; the smallest violation
(zero or two selected paths for some demand) costs $1$.

\paragraph{Latency objective.}
The raw cost of routing demand $k$ on path $p$ is bandwidth-weighted by
default,
\begin{equation}
  L_{k,p} \;=\; b_k \sum_{e \in p} \ell_e
  \qquad \text{(optionally unweighted: } L_{k,p} = \textstyle\sum_{e\in p}\ell_e\text{)},
\end{equation}
so that a 10\,Gbps demand on a slow path costs ten times a 1\,Gbps demand.
Latencies are normalized so that the spread of the latency term over the
feasible (one-hot) space equals a single tunable scale $c_{\mathrm{scale}}$:
\begin{equation}
  \label{eq:latnorm}
  \widetilde{H}_{\lat} = s_{\lat} \sum_{k,p} L_{k,p}\, x_{k,p},
  \qquad
  s_{\lat} = \frac{c_{\mathrm{scale}}}{\max(L^{\max} - L^{\min},\, \varepsilon)},
\end{equation}
where $L^{\min} = \sum_k \min_{p} L_{k,p}$ and
$L^{\max} = \sum_k \max_{p} L_{k,p}$ are the exact feasible-space extremes
(the latency term is separable per demand), and $\varepsilon$ guards a
degenerate zero spread.

\paragraph{Congestion (soft capacity).}
The utilization of link $e$ is linear in $x$,
\begin{equation}
  u_e(x) \;=\; \frac{1}{c_e} \sum_{(k,p)\,:\, e \in p} b_k\, x_{k,p},
\end{equation}
and the congestion term averages a convex profile $g$ over the $m_{\used}$
links that appear in at least one candidate path:
\begin{equation}
  \label{eq:cong}
  \widetilde{H}_{\cng} \;=\; \frac{c_{\mathrm{scale}}}{m_{\used}}
      \sum_{e \in E_{\used}} g\bigl(u_e(x)\bigr),
  \qquad
  g(u) =
  \begin{cases}
    u^{2} & \text{\texttt{quadratic} (default)},\\[2pt]
    \alpha u + \beta u^{2} & \text{\texttt{fortz-thorup-fit}}.
  \end{cases}
\end{equation}
No slack qubits and no hard capacity constraint are used: the quadratic
profile penalizes congestion smoothly and pushes solutions under capacity,
while hard feasibility is verified classically at decode time
(Section~\ref{sec:decode}). Since $u_e$ is dimensionless and $O(1)$ for sane
instances, the average keeps this term $O(c_{\mathrm{scale}})$.

The second profile calibrates $g$ to the canonical Fortz--Thorup link cost
$\Phi_e$ \cite{fortz2000}: the piecewise-linear convex function with
\begin{equation}
  \label{eq:phi}
  \Phi_e'(\ell) =
  \begin{cases}
     1    & \ell/c_e \in [0, 1/3)\\
     3    & \ell/c_e \in [1/3, 2/3)\\
     10   & \ell/c_e \in [2/3, 9/10)\\
     70   & \ell/c_e \in [9/10, 1)\\
     500  & \ell/c_e \in [1, 11/10)\\
     5000 & \ell/c_e \in [11/10, \infty)
  \end{cases}
  \qquad\Longrightarrow\qquad
  \Phi_e(c_e) = \tfrac{32}{3}\, c_e .
\end{equation}
A piecewise-linear function of a sum of binaries is not QUBO-expressible
without auxiliary variables, so the in-Hamiltonian surrogate is the
non-negative least-squares fit of $\alpha u + \beta u^2$ to $\Phi_e$ per unit
capacity on $u \in [0, 1.1]$. The hockey-stick tail of $\Phi$ drives the
unconstrained linear coefficient negative (which would \emph{reward} lightly
loaded links), so $\alpha$ is clamped at $0$; the fit is a tail-calibrated
pure quadratic, $\alpha = 0$, $\beta \approx 16.73$ --- about $16.7\times$ the
default congestion pressure. The linear part (when present) folds into the
linear coefficients of $H$; the quadratic part reuses the $u_e^2$ machinery.

\paragraph{Route changes (optional, default off).}
Given a current routing $r : k \mapsto p_{\mathrm{cur}}(k)$ for some or all
demands,
\begin{equation}
  H_{\switch} \;=\; \sum_{k \,\in\, \operatorname{dom} r}\;
      \sum_{p \neq p_{\mathrm{cur}}(k)} x_{k,p},
\end{equation}
so in the feasible subspace each rerouted demand costs exactly
$\lambda_{\switch}\, c_{\mathrm{scale}}$. The default is
$\lambda_{\switch} = 0$ (the term is skipped entirely).

\subsection{QUBO form}

Using $x^2 = x$ for binaries, each squared affine term in
\eqref{eq:onehot}--\eqref{eq:cong} expands to the standard QUBO monomials; for
a link $e$ with coefficients $a_{e,i} = b_{k(i)}/c_e$ over the paths crossing
it,
\begin{equation}
  u_e^2 \;=\; \sum_i a_{e,i}^2\, x_i \;+\; 2 \sum_{i<j} a_{e,i}\, a_{e,j}\, x_i x_j ,
\end{equation}
i.e.\ pairwise couplings between any two candidate paths sharing a link.
Cross-terms between two paths of the \emph{same} demand also appear; they are
inert on the one-hot-feasible subspace (at most one of the two variables is
$1$) and mildly reinforce $H_{\onehot}$ outside it, so they are kept. The
total Hamiltonian is therefore degree~2 --- an Ising model directly
implementable as a QAOA phase separator.

\subsection{Weights and normalization}
\label{sec:weights}

\begin{table}[ht]
\centering
\begin{tabular}{lll}
\toprule
Parameter & Default & Role \\
\midrule
$\lambda_{\onehot}$        & $4.0$ & one-hot penalty weight \\
$\lambda_{\cng}$          & $1.0$ & congestion weight \\
$\lambda_{\switch}$        & $0.0$ & route-change weight (off) \\
$c_{\mathrm{scale}}$       & $1.0$ & overall objective scale \\
congestion profile         & \texttt{quadratic} & $g(u)$ in \eqref{eq:cong} \\
\bottomrule
\end{tabular}
\caption{Default Hamiltonian weights.}
\label{tab:weights}
\end{table}

The normalization \eqref{eq:latnorm} fixes the feasible-space spread of
$\widetilde{H}_{\lat}$ to $c_{\mathrm{scale}}$ and keeps
$\widetilde{H}_{\cng} = O(c_{\mathrm{scale}})$, so the maximum feasible
objective gap is $\approx 2\,c_{\mathrm{scale}}$ and
$\lambda_{\onehot} = 4$ dominates it: breaking the one-hot constraint is
never energetically favorable. With $K$ demands the total spread of $H$ stays
$O(10)$, keeping $e^{-i\gamma H}$ away from $2\pi$ phase wrap-around for
$\gamma \in (0,\pi)$ --- the normalization strategy of Classiq's
network-traffic-optimization reference notebook \cite{classiqtelecom}.

\section{QAOA solution}
\label{sec:qaoa}

\subsection{Ansatz}

The circuit is the standard depth-$L$ QAOA \cite{farhi2014} over $N$ qubits,
built on Classiq as an alternation of an arithmetic phase gadget and a
transverse mixer:
\begin{equation}
  \label{eq:ansatz}
  \left| \psi(\gamma, \beta) \right\rangle
  \;=\;
  \prod_{l=1}^{L} \Bigl[ U_M(\beta_l)\; U_C(\gamma_l) \Bigr]\;
  \left| + \right\rangle^{\otimes N},
\end{equation}
\begin{equation}
  U_C(\gamma)\,|x\rangle = e^{-i \gamma H(x)}\,|x\rangle,
  \qquad
  U_M(\beta) = \bigotimes_{i=1}^{N} R_X^{(i)}(\beta),
  \quad R_X(\beta) = e^{-i \beta X / 2}.
\end{equation}
$U_C$ is generated directly from the polynomial $H(x)$ of
Section~\ref{sec:qubo} by Classiq's \texttt{phase} statement (the synthesis
engine compiles the degree-2 polynomial into single- and two-qubit phase
rotations); sign and half-angle conventions relative to \cite{farhi2014} are
reparameterizations absorbed by the classical optimizer. The identical Python
cost function is reused classically to score sampled bitstrings --- one
definition, two uses.

\subsection{Parameter optimization}

The variational objective is the shot-estimated energy
\begin{equation}
  F(\gamma, \beta) \;=\;
  \bigl\langle \psi(\gamma,\beta) \bigr| H \bigl| \psi(\gamma,\beta) \bigr\rangle
  \;\approx\; \sum_{x} \frac{n_x}{S}\, H(x),
\end{equation}
with $n_x$ the observed counts over $S$ shots. $F$ is minimized with COBYLA,
starting from the midpoint linear schedule
\begin{equation}
  \gamma_l = \pi\,\frac{2l-1}{2L}, \qquad
  \beta_l = \gamma_{L+1-l}, \qquad l = 1,\dots,L
\end{equation}
(ascending $\gamma$, descending $\beta$ --- an annealing-like ramp). Defaults:
$L = 3$ layers, $S = 2048$ shots, at most $60$ COBYLA iterations, fixed
simulator seed for reproducibility. After convergence the circuit is sampled
once more at the optimal parameters to obtain the solution distribution.

\subsection{Decoding, feasibility, and KPIs}
\label{sec:decode}

Each sampled bitstring is decoded classically:
\begin{itemize}
  \item \textbf{Routes:} demand $k$ is \emph{satisfied} iff exactly one
        $x_{k,p} = 1$; the selected paths determine per-link loads
        (unsatisfied demands carry no traffic).
  \item \textbf{Feasibility:} one-hot satisfaction per demand and hard
        capacity checks per link (\emph{load} $\le c_e$) --- the soft
        congestion term does not replace validation.
  \item \textbf{KPIs (raw units, never normalized):} per-link
        load and utilization, maximum utilization, capacity violations,
        total and bandwidth-weighted latency, demands satisfied, and the
        Fortz--Thorup congestion measure
        \begin{equation}
          \Phi = \sum_{e \in E} \Phi_e(\mathrm{load}_e),
          \qquad
          \Phi^{*} = \frac{\Phi}{\Phi_{\Uncap}},
          \qquad
          \Phi_{\Uncap} = \sum_{k} b_k \cdot \operatorname{hopdist}(s_k, t_k),
        \end{equation}
        with $\Phi_e$ the exact piecewise cost \eqref{eq:phi} (not the fit)
        and $\operatorname{hopdist}$ the unit-weight shortest-path distance.
        $\Phi^{*} = 1$ is the uncapacitated ideal; values from $1$ to roughly
        $\tfrac{32}{3} \approx 10.67$ mean all traffic flows within capacity;
        $\Phi^{*} \ge \tfrac{32}{3}$ signals a congested network
        \cite{fortz2000}.
\end{itemize}
For small instances an exhaustive scan of the $\prod_k |P_k|$ one-hot
assignments provides the exact optimum as a reference; QAOA quality is judged
by (i) convergence of $F$ toward that reference, (ii) the probability mass on
one-hot-feasible bitstrings, and (iii) the best sampled feasible solution
matching (or nearing) the reference optimum.

\section{Worked example (9 qubits)}
\label{sec:example}

\subsection{Instance}

Six routers \textsf{A}--\textsf{F}, nine directed links
(Table~\ref{tab:links}), three demands (Table~\ref{tab:demands}). Candidate
paths are the $\le 3$ latency-shortest simple paths per demand --- the
stand-in for pipeline step~2 --- giving $|P_1| = |P_2| = |P_3| = 3$ and
$N = 9$ qubits (Table~\ref{tab:indexmap}).

\begin{table}[ht]
\centering
\begin{minipage}[t]{0.48\textwidth}
\centering
\begin{tabular}{lcc}
\toprule
Link & $c_e$ & $\ell_e$ \\
\midrule
A$\to$B & 10 & 1 \\
B$\to$D & 5  & 1 \\
A$\to$C & 10 & 2 \\
C$\to$D & 10 & 2 \\
B$\to$C & 5  & 1 \\
D$\to$E & 10 & 1 \\
D$\to$F & 5  & 2 \\
E$\to$F & 10 & 1 \\
C$\to$E & 10 & 3 \\
\bottomrule
\end{tabular}
\caption{Links (capacity, latency).}
\label{tab:links}
\end{minipage}\hfill
\begin{minipage}[t]{0.48\textwidth}
\centering
\begin{tabular}{lccc}
\toprule
Demand & $s_k$ & $t_k$ & $b_k$ \\
\midrule
$d_1$ & A & F & 4 \\
$d_2$ & A & E & 3 \\
$d_3$ & B & F & 3 \\
\bottomrule
\end{tabular}
\caption{Traffic demands.}
\label{tab:demands}
\end{minipage}
\end{table}

\begin{table}[ht]
\centering
\begin{tabular}{clclcc}
\toprule
Qubit & Demand & $p$ & Path & $\sum_{e\in p}\ell_e$ & $L_{k,p} = b_k \ell(p)$ \\
\midrule
$x_0$ & $d_1$ & 0 & A--B--D--F     & 4 & 16 \\
$x_1$ & $d_1$ & 1 & A--B--D--E--F  & 4 & 16 \\
$x_2$ & $d_1$ & 2 & A--C--D--F     & 6 & 24 \\
$x_3$ & $d_2$ & 0 & A--B--D--E     & 3 & 9  \\
$x_4$ & $d_2$ & 1 & A--C--E        & 5 & 15 \\
$x_5$ & $d_2$ & 2 & A--B--C--E     & 5 & 15 \\
$x_6$ & $d_3$ & 0 & B--D--F        & 3 & 9  \\
$x_7$ & $d_3$ & 1 & B--D--E--F     & 3 & 9  \\
$x_8$ & $d_3$ & 2 & B--C--D--F     & 5 & 15 \\
\bottomrule
\end{tabular}
\caption{Candidate paths and the flat index map ($o_1 = 0$, $o_2 = 3$,
$o_3 = 6$).}
\label{tab:indexmap}
\end{table}

The normalization constants evaluate to $L^{\min} = 34$, $L^{\max} = 54$,
hence $s_{\lat} = 1/20 = 0.05$; all $m_{\used} = 9$ links are crossed by some
candidate path, so the congestion prefactor is
$\lambda_{\cng} c_{\mathrm{scale}} / m_{\used} = 1/9$; the scaled one-hot
weight is $4$. The uncapacitated Fortz--Thorup bound is
$\Phi_{\Uncap} = 4\cdot3 + 3\cdot2 + 3\cdot2 = 24$.

\subsection{Reference solution}

Table~\ref{tab:results} reports the exact optimum of $H$ over all $27$
one-hot assignments (the reference QAOA must reproduce), for three
configurations of the congestion term. KPIs are computed by the decoder in
raw units.

\begin{table}[ht]
\centering
\small
\begin{tabular}{lccccc}
\toprule
Congestion term & $(p_{d_1}, p_{d_2}, p_{d_3})$ & $\max_e u_e$ &
Cap.\ viol. & $\sum_k b_k \ell(p_k)$ & $\Phi^{*}$ \\
\midrule
off ($\lambda_{\cng}=0$)
  & (A--B--D--F,\; A--B--D--E,\; B--D--F) & 2.00 & 2 & 34 & 1276.1 \\
\texttt{quadratic} (default)
  & (A--B--D--E--F,\; A--B--D--E,\; B--D--F) & 2.00 & 1 & 34 & 952.0 \\
\texttt{fortz-thorup-fit}
  & (A--C--D--F,\; A--C--E,\; B--D--E--F) & \textbf{0.80} & \textbf{0} & 48 & \textbf{2.08} \\
\bottomrule
\end{tabular}
\caption{Exact optimum of $H$ under three congestion settings (all other
weights at their defaults, Table~\ref{tab:weights}). Latency is
bandwidth-weighted; $\Phi^{*}$ uses the exact piecewise link cost.}
\label{tab:results}
\end{table}

The example is built so the pure-latency optimum overloads shared links: all
three demands prefer the cheap corridor through B$\to$D (capacity~5), which
with $\lambda_{\cng} = 0$ carries load $10$ ($u = 2.0$) alongside a second
violation on D$\to$F. Switching on the default quadratic term changes the
optimum (it reroutes $d_1$ off D$\to$F, halving the Fortz--Thorup cost) ---
the congestion term demonstrably reshapes the routing --- but at
$\lambda_{\cng} = 1$ the latency term still dominates and B$\to$D stays
overloaded. The Fortz--Thorup-calibrated profile ($\beta \approx 16.73$)
prices overload realistically and drives the optimum to a fully feasible
routing: zero capacity violations, maximum utilization $0.80$, and
$\Phi^{*} = 2.08$ (well below the congestion threshold $32/3$), bought with
$+41\%$ bandwidth-weighted latency ($34 \to 48$). This is exactly the
congestion-relief-versus-latency trade-off of the challenge KPIs, exposed
through the $\lambda_{\cng}$ / profile knobs.

\subsection{QAOA run}

The demo notebook (\texttt{qaoa\_routing.ipynb}) synthesizes the 9-qubit,
$L = 3$ ansatz \eqref{eq:ansatz} for this instance, optimizes
$(\gamma, \beta)$ with COBYLA ($S = 2048$ shots, $\le 60$ iterations) on the
Classiq executor, and reports: the convergence of $F$ against the
brute-force optimum of Table~\ref{tab:results} as a reference line, the
probability mass on one-hot-feasible bitstrings, the top-10 sampled solutions
with full KPI rows, and the best feasible sample versus the exact optimum.
Success is judged by the soft criteria of Section~\ref{sec:decode} (COBYLA on
shot noise is not asserted to hit the exact optimum).

\section{Scaling and outlook}
\label{sec:outlook}

\paragraph{Qubit economy.} The path-based encoding uses
$N = \sum_k |P_k| \approx K \bar{k}$ qubits for $K$ demands with $\bar{k}$
candidate paths each, independent of network size; an edge-based
flow-conservation encoding of the same 6-node instance
\cite{classiqtelecom} needs $K \cdot |E| = 27$ qubits versus our $9$, and the
gap widens with the network. The price is the classical restriction to
candidate sets --- path quality is bounded by step~2 of the pipeline.

\paragraph{Knobs.} $\lambda_{\cng}$ and the congestion profile trade latency
against congestion relief (Table~\ref{tab:results}); $\lambda_{\switch}$ with
a current routing penalizes unnecessary route changes;
$c_{\mathrm{scale}}$, layer count $L$, shots, and iteration budget control the
QAOA itself.

\paragraph{Next steps (pipeline step 5).} Compare QAOA-selected routings
against a classical baseline (MILP over the same candidate sets, and the
Fortz--Thorup LP bound for unrestricted routing) on the full KPI set:
most-congested-link relief, average latency, demand satisfaction, route
changes, and zero capacity/resiliency violations.

\begin{thebibliography}{9}

\bibitem{farhi2014}
E.~Farhi, J.~Goldstone, and S.~Gutmann,
\emph{A Quantum Approximate Optimization Algorithm},
arXiv:1411.4028 (2014).

\bibitem{fortz2000}
B.~Fortz and M.~Thorup,
\emph{Internet Traffic Engineering by Optimizing OSPF Weights},
Proc.\ IEEE INFOCOM 2000, pp.~519--528.

\bibitem{classiqtelecom}
Classiq Technologies,
\emph{Network Traffic Optimization} (QAOA application notebook),
\url{https://github.com/Classiq/classiq-library/tree/main/applications/telecom/network_traffic_optimization}.

\end{thebibliography}

\end{document}
